\section{The elliptic suspension and its projective example}
\label{sec:elliptic-series}

Let $f$ be a smooth symplectomorphism of standard $\RR^4$.
Take the quotient of $\RR^4\times\RR^2$ by
\[
  (x,s,t)\longmapsto(f(x),s+1,t),\qquad
  (x,s,t)\longmapsto(x,s,t+1).
\]
Give the horizontal planes the complex structure of the base elliptic
curve and the vertical planes any compatible almost-complex structure.
Section~\ref{sec:suspension} shows that every compact irreducible
pseudoholomorphic curve is the horizontal leaf of a finite orbit.
Let $U\Subset\RR^4$ be open. Assume that any finite orbit contained
in $\bar U$ lies in $U$. For $k\ge1$ set
\[
 U_k=\bigcap_{j=0}^{k-1}f^{-j}U,\qquad
 I_k=\deg(\operatorname{id}-f^k,U_k,0).
\]
These are the local fixed point indices of the iterates, restricted to
orbits staying in $U$. The integrality of their M\"obius transforms is
the subject of Dold's iteration congruences \cite{Dold}.

Let $H$ denote the Chern-number-zero, virtual-dimension-zero part of
Pardon's rational algebra of semi-Fano complex one-cycles \cite{pardon}.
Here semi-Fano means that every irreducible component has nonnegative
Chern number. All components at the complex targets below have Chern
number zero.
Use the comparison $\Phi:H\to S$ to the algebra defined using smooth
almost-complex families \cite[\S2]{WueProduct}. A compact open set of cycles in
a target defines an element of $S$ by its characteristic function. Take
the series $P_U(q)$ of these classes for the cycles supported on orbits
in $U$, graded by projection degree over the elliptic base, with constant
term one. The compactness needed for this definition holds separately
in each degree. For a nontorsion degree-zero line $L$ on an elliptic curve
$C$, let $F_1(q)$ be the series in $H$ of the point insertions at
$m[C]$ in the total space of $L\oplus L^{-1}$, including the empty cycle.
These are the scalar local classes in genus one
\cite[Definition~4.7]{pardon}. Put $E(q)=\Phi(F_1(q))$.

\begin{theorem}\label{thm:index}
With the isolation assumption above,
\[
 P_U(q)=\prod_{d\ge1} E(q^d)^{b_d},\qquad
 b_d=\frac{1}{d}\sum_{k\mid d}\mu(d/k) I_k\in\ZZ .
\]
Equality of two such series is equivalent to equality of all their local
iterate indices. Through degree $K$ the formula requires isolation only
for orbits of period at most $K$.
\end{theorem}

Ionel and Parker related the numerical Gromov series of symplectic
mapping tori to the Lefschetz zeta function of the return map
\cite[Theorems~0.1 and~4.1]{IPmaps}. In real dimension six their formula
uses the genus-one numerical invariants. Dold's indices already determine
the product expansion of the local dynamical zeta function.
Theorem~\ref{thm:index} identifies their contribution to the
cycle classes themselves, before an enumerative evaluation. The step
that makes this identification possible is the calculation at one orbit.
A saddle-center bifurcation expresses the positive hyperbolic contribution
as the inverse of the elliptic one. At period doubling an additional
embedded torus accounts for a factor in twice the primitive class.
Ionel and Parker describe the corresponding relations for numerical
generating functions and their connection with Taubes's torus counts
\cite[Remark~3.4]{IPmaps}. Here compact continuation gives the relations
for the classes represented by the characteristic functions of the entire cycle loci. They determine
the contribution of every four-dimensional symplectic return whose
relevant iterates are nondegenerate. For each degree bound a small
Hamiltonian perturbation then treats arbitrary degeneracies in the isolated set.

\subsection{Compact curves in a suspension}\label{sec:suspension}

The product horizontal distribution on the covering space descends to the
suspension. A metric on its vertical symplectic bundle gives a compatible
vertical complex structure by polar decomposition. Choices of compatible
structures form a contractible space and can be made smoothly along a
path of symplectomorphisms. Use the direct sum of the vertical and
horizontal structures throughout. The standard symplectic form of the
fiber descends to a closed two-form $\Omega$ whose kernel is the horizontal
plane. It is pulled back from the five-dimensional mapping torus over the
$s$ circle. That mapping torus has the homotopy type of a circle, since
its fiber is contractible. In particular $\Omega$ is exact.

\begin{lemma}\label{lem:sheets}
The compact irreducible pseudoholomorphic curves are the suspended
periodic orbits. A curve from an orbit of least period $j$ is an embedded
elliptic curve of degree $j$ over the base and has Chern number zero.
\end{lemma}
\begin{proof}
On the normalization of a compact irreducible curve the integral of
$\Omega$ is zero. In a holomorphic local coordinate its area density is
the squared norm of the vertical part of the first derivative, using the
compatible vertical metric. The vertical derivative vanishes everywhere.
The normalization therefore maps to a single horizontal leaf, where
projection to the base is locally biholomorphic. A nonconstant map from
a compact connected curve to that leaf is open and has closed image, by
lifting it in the flat foliation charts. Its image is the whole leaf.
A compact leaf meets a transverse fiber in one finite orbit, and
conversely a finite orbit of length $j$ gives the embedded quotient with
periods $j,i$ in the base coordinate $s+it$. The tangent line to the
leaf is trivial. Its vertical normal symplectic bundle is pulled back
from the $s$ circle, so its first Chern class vanishes.
\end{proof}

In a compact-fiber suspension the same proof identifies all the curves
of zero $\Omega$-area, since the area density stays nonnegative.
For a cycle in the local suspension, intersect with the fiber over the
origin of the elliptic base. This gives an $f$-invariant multiset of
length equal to the degree of the cycle. Conversely, an invariant
multiset is a union of finite orbits with constant multiplicity on
each orbit, so determines a unique cycle. This correspondence allows
several formerly distinct orbits to specialize to the same multiple
leaf. In a flat target chart the invariant points give constant
horizontal sheets. Convergence of the finite multiset thus gives
Hausdorff convergence of the supports and weak convergence of their
currents, the topology used in \cite[\S2]{WueProduct}. Finitely many strips
cover the base and the same argument works when their gluing map varies.

For a fixed degree $m$ consider all the invariant multisets in
$\operatorname{Sym}^m\bar U$. They form a closed subset of this compact
space. Every constituent orbit has period at most $m$, so boundary
isolation puts the entire multiset inside $U$. Its continuous
parametrization of the cycles is a homeomorphism onto its image by
compactness and the Hausdorff property. This image is open in the
degree-$m$ cycle space. Indeed, for a retained cycle all its intersection
points with the transverse fiber lie in $U$. Nearby supports avoid the
closed complement on that fiber, by compactness of the original
support and upper Hausdorff continuity. The degree itself is locally
constant, by testing the currents against the base area form with a
cutoff equal to one near the compact support. Over a compact parameter
interval with boundary isolation the argument applies in the relative
symmetric power and gives a compact open family. Its end characteristic
functions agree in $S$ by the interval relation \cite[\S2]{WueProduct}.

For disjoint neighborhoods of finitely many closed leaf supports, choose
disjoint open tubes in the suspension containing those supports.
For each distribution of the total degree among the tubes, external
product followed by open extension represents the characteristic function
of the corresponding cycles. Summing over all distributions gives the
product of the local series. A single orbit of least period $j$ has
deck gluing $f^j$ in its tubular germ over the elliptic curve of first
period $j$. In terms of the original base degree its series is obtained
from the local one by the substitution $q\mapsto q^j$.

\subsection{The contribution of a nondegenerate orbit}\label{sec:returns}

Let $x$ have least period $d$ under the symplectomorphism $f$, and put
$T=D_xf^d$. A transverse disc at $x$, followed for $d$ periods,
gives the germ of the primitive leaf as a suspension of $f^d$ over an
elliptic curve with first period $d$. We use $q$ for the class of this
leaf. In the grading by degree over the original base its contribution
is obtained by substituting $q^d$ for $q$. Our deck transformation
advances the first base coordinate by one and acts on the fiber by
$f$. Thus positive parallel transport has holonomy $f^{-1}$.
The periodic orbits are the same, and
$\det(I-T^{-j})=\det(I-T^j)$ since $T$ is symplectic and its real
dimension is even.

\begin{proposition}\label{prop:return}
Suppose $\det(I-T^j)\ne0$ for $1\le j\le K$, and put
$\sigma_j=\operatorname{sign}\det(I-T^j)$ in this range. For $K\ge2$
the local series of the primitive leaf is
\begin{equation}\label{eq:primitive-return}
 E(q)^{\sigma_1} E(q^2)^{(\sigma_2-\sigma_1)/2}
 \pmod {q^{K+1}} .
\end{equation}
The sign $\sigma_j$ is $\sigma_1$ for odd $j$ and $\sigma_2$ for
even $j$. For $K=1$ the series is $E(q)^{\sigma_1}$ modulo $q^2$.
If every iterate is nondegenerate, \eqref{eq:primitive-return} holds
as an identity of formal series.
\end{proposition}

For a sum of two symplectic planes write $P_{A,B}$ for its local
series, using $e,+,-$ for elliptic, positive hyperbolic and negative
hyperbolic blocks. Fix the second block $B$ with its first $K$ iterates
nondegenerate. A saddle-center bifurcation and a period doubling give,
respectively, modulo $q^{K+1}$,
\begin{equation}\label{eq:return-walls}
 P_{+,B}(q)P_{e,B}(q)=1,\qquad
 P_{-,B}(q)P_{e,B^2}(q^2)=P_{e,B}(q).
\end{equation}
Each elliptic symbol refers to the rotation at its indicated endpoint.
The two points exchanged at period doubling give one leaf of base degree
two. Its return on the second plane is $B^2$, which explains both the
single factor and the substitution in the second identity.

\begin{proof}
Work in a Darboux chart centered at the return point and denote
the primitive return germ by $g$. For $0<\lambda\le1$ the maps
$g_\lambda(y)=\lambda^{-1}g(\lambda y)$ are symplectic and extend
smoothly to $g_0=T$. For each $j\le K$ their $j$th iterates have
the form $T^j y+O(\|y\|^2)$ uniformly in $\lambda$. Since the
finitely many $I-T^j$ are invertible, there is a common ball in
which zero is the only fixed point of these iterates. To realize the
family of germs in a suspension, the isotopy $T^{-1}g_\lambda$
has a Hamiltonian on a smaller ball, obtained by integrating the
closed one-form of its generating vector field. A cutoff of this
Hamiltonian, followed by composition with $T$, extends the family
to $\RR^4$ while preserving it on a neighborhood of zero. The
characteristic insertion of each multiple through degree $K$ has
compact support along the central leaf over the parameter interval.
Continuation identifies its class with that of the linear suspension.
The same argument identifies the insertions along any path of
linear maps satisfying the stated nondegeneracy conditions.

First suppose $T$ is a sum of two real symplectic planes. A planar
block is elliptic if it is conjugate to a rotation, and positive or
negative hyperbolic if its eigenvalues are distinct positive or
negative reciprocal reals, respectively. Blocks of
either hyperbolic type can be joined within that type by symplectic
conjugation and variation of the real eigenvalues. Such paths avoid
all resonances. We retain the angle of each elliptic block until its
contribution has been determined.

For two elliptic blocks choose a complex structure making both
rotations unitary. The suspension is the total space of two flat
holomorphic lines on the elliptic base. Each power of either line
of order $1,\ldots,K$ is nontrivial, since a holomorphically trivial
unitary flat line has trivial holonomy. For a cycle of degree
$m\le K$, the elementary symmetric coefficients of its fiber
projection to a line $L$ are holomorphic sections of $L^j$ for
$1\le j\le m$. This follows across branch points from the
characteristic polynomial on the finite direct image of the
normalization, using the cycle multiplicities. All these coefficients
vanish. Thus the degree-$m$ cycle is the $m$-fold zero section.
In the elliptic Picard variety the complement of the torsion points
of orders at most $K$ is path connected. Move the two lines there
independently to a nontorsion pair $L,L^{-1}$. At every stage the
same multiple zero sections give compact point insertions. At the inverse pair these are precisely the point insertions defining
$F_1$. It follows that $P_{e,e}=E$, for either angular
orientation and every base modulus
\cite[Definition~4.7]{pardon}, \cite[\S3]{WueProduct}.

Fix a second planar block $B$. On the first plane consider the
time-$\delta$ Hamiltonian map of
\[
 H_a(u,v)=\tfrac12v^2+\tfrac13u^3-au .
\]
For small $a<0$ it has no equilibrium in a disc about the origin.
For $a>0$ the critical points at $u=\sqrt a,-\sqrt a$, $v=0$,
give an elliptic and a positive hyperbolic fixed point. For a
second family take negation composed with the time-$\delta$ map of
\[
 J_a(u,v)=\tfrac12v^2+\tfrac14u^4-\tfrac a2u^2 .
\]
Negation commutes with this flow. At $a<0$ the origin is elliptic.
At $a>0$ it is negative hyperbolic and the two minima at
$u=\pm\sqrt a$ are exchanged. They give one orbit of period two.
The second iterate at a minimum is the time-$2\delta$ elliptic
linearized flow on the first plane and acts by $B^2$ on the second.

For both families take a Hamiltonian cutoff equal to one on a
larger disc, even under negation in the second family. The resulting
vector fields have a common Lipschitz bound $L_0$ for $a$ in a
compact interval about zero. A nonconstant periodic trajectory of
period $\tau$ has $\tau L_0\ge1$. Indeed, its velocity has mean
zero, and the difference of the velocities at two times is bounded
by $L_0\tau$ times the maximum speed. Evaluate the mean difference
at a point of maximum speed. Choose $2K\delta L_0<1$. In the first
family any point fixed by the $j$th iterate is fixed by the
time-$j\delta$ flow. In the second it is fixed by the time-$2j\delta$
flow. Hence the only periodic points through the cutoff in the
smaller disc are the equilibria just described, including the
degenerate one at $a=0$. The second block contributes only its
origin. Choose the regular endpoint parameters to avoid the finitely
many forbidden elliptic resonances. All retained orbits lie in the
interior of the disc throughout. Their invariant multisets in each
degree form a compact family, so the interval relation applies to
the whole insertion across either collision.

Excision into the distinct orbit neighborhoods at the endpoints now gives
\eqref{eq:return-walls}. In the second identity the squared block is
needed only through degree $\lfloor K/2\rfloor$. At $K=1$ its factor
is omitted.

First use \eqref{eq:return-walls} with $B$ elliptic. The
two-elliptic calculation at each endpoint gives $P_{+,e}=E^{-1}$
and $P_{-,e}=E/E(q^2)$ for every admissible second angle.
Interchanging the planes now determines $P_{e,B}$ for either
hyperbolic type and every admissible first angle. Applying
\eqref{eq:return-walls} once more yields
\[
 P_{+,+}=E,\qquad P_{+,-}=E(q^2)/E,\qquad P_{-,-}=E .
\]
For a planar elliptic block the determinant signs are positive
at every allowed power. A positive hyperbolic block has negative
sign at every power, and a negative hyperbolic block has positive
sign at odd powers and negative sign at even powers. Multiplying
the two signs proves the proposition for the split cases.

For general $T$, take a small coordinate ball in
$\operatorname{Sp}(4,\RR)$ on which the first $K$ determinants
remain nonzero. This ball contains a matrix of simple spectrum.
Indeed the discriminant of the characteristic polynomial is a real
analytic function on the connected symplectic group and is nonzero
on a diagonal matrix with distinct reciprocal pairs. A path in the
ball to a simple-spectrum matrix preserves the local insertion and
every relevant sign. In the complexification, the symplectic form
pairs eigenlines only when their eigenvalues are reciprocal. A real
reciprocal pair or a conjugate pair on the unit circle therefore
determines a real symplectic plane, reducing these cases to the
calculation above.

The remaining spectrum consists of a nonreal quartet
$\lambda,\bar\lambda,\lambda^{-1},\bar\lambda^{-1}$ off the unit
circle. The real stable and unstable planes are Lagrangian by
invariance of the symplectic form under forward and backward
iteration, respectively. In dual bases of the two planes the
matrix has the form $\operatorname{diag}(A,A^{-\mathsf T})$.
Choose the stable basis to give $A=rR_\theta$, with $0<r<1$
and $R_\theta$ a rotation, and change the unstable basis dually.
Deforming $\theta$ to zero keeps the spectrum off the unit circle.
The endpoint splits into two positive hyperbolic planes with
eigenvalues $r,r^{-1}$ on each. Its contribution is $E$, and all
its determinant signs are positive. This proves
\eqref{eq:primitive-return} through the prescribed cutoff for
every $T$. Open excision identifies the coefficients obtained in
smaller neighborhoods as $K$ increases, giving the formal identity
when all iterates are nondegenerate.
\end{proof}

\subsection{Fixed point indices of the iterates}\label{sec:indices}

Fix $K\ge1$ and assume that every orbit of period at most $K$
contained in $\bar U$ lies in $U$. This condition persists under a
sufficiently small perturbation of $f$ in the compact-open topology.
Indeed, from a sequence of perturbed orbits in $\bar U$ meeting
$\partial U$, take a subsequence of one fixed period and choose a
boundary point on each orbit. Compactness gives a limiting orbit of
$f$ in $\bar U$ with a point on $\partial U$.

For a path $f_t$ satisfying this condition, put
\[
 U_{k,t}=\bigcap_{i=0}^{k-1}f_t^{-i}U,\qquad 1\le k\le K .
\]
A zero of $\operatorname{id}-f_t^k$ on $\partial U_{k,t}$ would give
a forbidden boundary orbit. The retained zero sets have compact
union in the product of the parameter interval with $\bar U$.
Indeed, one may replace $U$ by $\bar U$ in their definition and
then use the isolation assumption. At any $t_0$ choose an open
set $W\Subset U_{k,t_0}$ containing the retained zeros. For nearby
$t$ we have $\bar W\subset U_{k,t}$ by continuity, and $W$ still
contains every retained zero by compactness. If there is no zero at
$t_0$, this remains true on a short interval. Excision and homotopy
invariance of the Brouwer degree on $W$ show that $I_k$ is constant
along the path. The series $P_U$ is constant through degree $K$ by
the compact family of invariant multisets in
\S\ref{sec:suspension}.

\begin{lemma}\label{lem:finite-perturbation}
There is an arbitrarily small Hamiltonian perturbation $g$ of $f$
for which every orbit in $U$ of least period $j\le K$ is
nondegenerate at every iterate of order $aj\le K$.
\end{lemma}
\begin{proof}
We successively postcompose the map with compactly supported
Hamiltonian diffeomorphisms, proceeding in order of least period
$j=1,\ldots,K$. Suppose that the orbits of smaller least period
already have the required property. There are finitely many of
them, by compactness and nondegeneracy. Around their points choose
disjoint inverse function neighborhoods for the finitely many
iterate equations of order at most $K$. For an iterate fixing a
point, its unique continued zero is the continuation of the
primitive periodic point. For every other iterate the
neighborhood can be chosen disjoint from its fixed points.
These properties persist under small $C^1$ perturbations.

Write $F$ for the current map. Outside smaller neighborhoods of the
reserved points, the set of $x\in\bar U$ with $F^j x=x$ and orbit
in $\bar U$ is compact and consists of points of exact period $j$.
At each such $x$ take a Darboux ball disjoint from the other points
of its orbit. Four smooth Hamiltonians supported in this ball can
be chosen with Hamiltonian vectors spanning $T_x\RR^4$ at $x$.
For the postcomposition of $F$ with their flows, the derivatives
of the $j$th iterate in these four time parameters receive a
contribution only at the last return to $x$. They therefore span
$T_x\RR^4$. This rank condition persists on a neighborhood of $x$.

Finitely many such neighborhoods cover the compact set. Compose all
their flows in one family with independent parameters. At zero
parameters the four columns belonging to each chosen ball retain
the rank just computed, even if different balls overlap. The
parameter-dependent iterate is consequently transverse to the
diagonal near all the relevant zeros. Compactness and boundary
isolation place all nearby retained zeros in these neighborhoods or
in the reserved ones. By parametric transversality an arbitrarily
small choice gives transversality at each of the remaining fixed
points of the $j$th iterate. They now form a finite set.

Select one point on each of the new orbits and choose disjoint
Darboux balls there that avoid every other orbit point of period
at most $j$ in $U$. A quadratic Hamiltonian centered at the chosen
point and cut off in its ball fixes the orbit and multiplies the
linear return on the left by a prescribed symplectic matrix near
the identity: take the Hamiltonian logarithm of this matrix as the
infinitesimal linear map. For each $1\le a\le\lfloor K/j\rfloor$
the equation $\det(1-T^a)=0$ defines a proper real analytic subset
of the connected symplectic group. We can therefore perturb the
primitive return $T$ arbitrarily little so that all these
determinants are nonzero. The other points of the orbit have
conjugate return matrices. The previously transverse fixed points
of the iterates of order at most $j$ persist uniquely, and
compactness excludes new ones for a sufficiently small perturbation.
This completes the induction. All the composing flows can be taken
small enough to give a path from $f$ to $g$ preserving boundary
isolation through $K$.
\end{proof}

\begin{proof}[Proof of Theorem~\ref{thm:index}]
Use the perturbation of Lemma~\ref{lem:finite-perturbation}. For
an orbit of $g$ of least period $j\le K$, let $T$ be its primitive
return matrix and $s=\operatorname{sign}\det(1-T)$. If $2j\le K$,
put $t=\operatorname{sign}\det(1-T^2)$. The sign at the $a$th
power of $T$ is $s$ for odd $a$ and $t$ for even $a$, as follows
from the planar blocks and the symplectic deformations in
Proposition~\ref{prop:return}. The $j$ points of the orbit thus
contribute $js$ to $I_k$ for odd $k/j$, $jt$ for even $k/j$,
and zero if $j$ does not divide $k$.

By Proposition~\ref{prop:return} the orbit contributes
\[
       E(q^j)^s\, E(q^{2j})^{(t-s)/2}
\]
through degree $K$, omitting the second factor when $2j>K$.
Excision into disjoint orbit tubes multiplies these factors.
Write $e_d$ for the resulting exponent of $E(q^d)$ for $d\le K$.
Each $e_d$ is an integer. In the divisor sum $\sum_{d\mid k}d e_d$
the first factor of this orbit contributes $js$ at each multiple
of $j$. Its second adds $j(t-s)$ at the even multiples. These
are exactly the local index contributions above. Hence
\[
        I_k=\sum_{d\mid k} d e_d,\qquad 1\le k\le K .
\]
M\"obius inversion gives the stated formula for $b_d=e_d$.
Both the series and the indices agree for $f$ and $g$ through
this order. The exponents thus depend only on the original
indices, and the products for successive $K$ agree in the
overlapping degrees.

Finally write $E(q)=1+cq+O(q^2)$. The class $c=\Phi(x_{1,1})$
is nonzero by multiplicity-one detection \cite[Theorem~5.1]{WueProduct}.
If two exponent sequences first differ at $d$, their products
first differ in degree $d$ by the difference of the two exponents
times $c$. This is nonzero over $\QQ$. Equality of the series is
therefore equivalent to equality of the exponents and hence, by
the divisor sum, to equality of the iterate indices.
\end{proof}

In particular, when the retained orbits split among disjoint open
sets, each orbit lying wholly in one of them, index excision gives
the product of the corresponding series. The moving-domain argument
also applies along any path with boundary isolation. It allows
permutation of individual orbits while preserving the series
graded by degree over the base.


Under an isolated deformation the retained periodic orbits can be
permuted. The index formula recovers exactly the unlabelled iterate
indices. In the following example the labels in integral homology
also distinguish the individual section orbits.

\subsection{A projective quotient with sixteen section classes}\label{sec:compact}

Let $A$ be the product of two elliptic curves with a translation-invariant
product K\"ahler form $\omega_A$. Set
\[
 E'=\CC/(2\ZZ+i\ZZ),\qquad B=\CC/(\ZZ+i\ZZ),\qquad
 X=(A\times E')/\langle (x,z)\longmapsto(-x,z+1)\rangle .
\]
The projection $p:X\to B$ is an abelian surface bundle. The involution
acts freely and preserves a nowhere-vanishing holomorphic three-form,
so $X$ is smooth and projective with $K_X$ trivial. Its invariant
one-forms come from $E'$, and its invariant two-forms come from $A$.
In particular $h^{1,0}(X)=h^{2,0}(X)=1$. The Hodge vanishing assumptions
of the compact product theorem of \cite{WueProduct} fail for this quotient.

Write $\Omega$ for the closed semipositive form on $X$ induced by
$\omega_A$. Let $\Gamma^0$ consist of the integral homology classes
of effective one-cycles on $X$ that pair to zero with $\Omega$,
including the class of the empty cycle. This is a face of the effective
monoid since $\Omega$ pairs nonnegatively with every irreducible curve.
For $\gamma\in\Gamma^0$ let $U_{X,\gamma}\in H$ be the class of
the characteristic function of the class-$\gamma$ locus in the cycle
space. With $\Lambda=H_1(A;\ZZ)$, identify $A[2]$ with
$\Lambda/2\Lambda$ by $v\mapsto v/2\bmod\Lambda$. Denote by
$\beta_v$ the class of the section through this point and by
$F_1(q)$ the scalar local series of an elliptic curve.

\begin{theorem}\label{thm:face}
The sixteen classes $\beta_v$ are distinct and generate $\Gamma^0$.
They satisfy $2\beta_v=2\beta_0$ for every $v$. In the monoid
algebra over $H$ completed by degree over $B$ one has
\begin{equation}\label{eq:face}
  \sum_{\gamma\in\Gamma^0}U_{X,\gamma}Q^\gamma
       =F_1(Q^{2\beta_0})^{-8}
          \prod_{v\in\Lambda/2\Lambda}F_1(Q^{\beta_v}).
\end{equation}
\end{theorem}

\begin{proof}
Write $z=s+it$. As a smooth manifold $X$ is $M\times S^1_t$,
where $M$ is the mapping torus of negation on $A$. On its universal
cover a lattice translation $T_v$ acts on the fiber and the first
period acts by $R(x,s)=(-x,s+1)$. The relation
$RT_vR^{-1}=T_{-v}$ gives
\[
 H_1(M;\ZZ)=(\Lambda/2\Lambda)\oplus\ZZ .
\]
The positively oriented $s$-loop based at $v/2$ lifts to the deck
transformation $T_vR$, of class $(v,1)$. Taking its cross product with
$[S^1_t]$ gives $\beta_v$ with its complex orientation. This cross
product identifies $H_1(M;\ZZ)$ with a direct summand of
$H_2(X;\ZZ)$ by the integral K\"unneth formula. It follows that the
sixteen labels are distinct and that their doubles coincide.

On the normalization of an irreducible curve the pullback of $\Omega$
is pointwise nonnegative. If its integral vanishes, the differential
has zero vertical component, so the curve is a compact horizontal
leaf, as in \S\ref{sec:suspension}. At the complex structure on $X$
these leaves correspond to the one- and two-element orbits of
negation. A two-element orbit traverses two $s$-periods, with deck
transformation $R^2$ and zero lattice translation. The class of its
leaf is therefore $2\beta_0$. This proves the assertion about the
monoid. Its generators all have base degree one. There are finitely
many labels of bounded base degree, and any ample degree on $X$
gives the same completion on this face since the degrees of its
finitely many generators have positive minimum and finite maximum.

\smallskip
\noindent\emph{Deformation of the return.}
Choose period-one real coordinates $(u_1,v_1,u_2,v_2)$ in which
$\omega_A$ is a positive constant multiple of $du_j\wedge dv_j$
on the $j$th plane. Let $\varphi^\tau$ be the Hamiltonian flow of
\[
 G(u_1,v_1,u_2,v_2)
       =\sum_{j=1}^2\bigl(\cos(2\pi u_j)+\cos(2\pi v_j)\bigr).
\]
It commutes with negation. Put $f_\tau=(-1)\circ\varphi^\tau$
and suspend $f_\tau$ as in \S\ref{sec:suspension}, with descended
vertical form $\Omega_\tau$. Choose compatible vertical complex
structures smoothly in the parameter, starting with the product
complex structure, and retain the flat horizontal complex line
throughout. Under the resulting identification of the underlying
manifolds, $[\Omega_\tau]=[\Omega]$. Indeed, choose a smooth
function $h$ on the fundamental strip $0\le s\le1$, identically zero
near $0$ and one near $1$. The map from the initial suspension to the
one at $\tau$ given on this strip by
$(x,s,t)\mapsto(\varphi^{\tau h(s)}x,s,t)$ respects the seams.
The change of vertical form is $\pm\tau h'(s)\,ds\wedge dG$.
It has a primitive $\mp\tau h'(s)G\,ds$ vanishing near both seams.
This identification fixes the sixteen section tori since the points
of $A[2]$ are equilibria.

For any parameter, every component in a cycle of a face class is
horizontal, by the nonnegativity of its vertical area. At a fixed
base degree $m$, intersecting with the fiber over $0\in B$ identifies
the horizontal cycles with the $f_\tau$-invariant elements of
$\operatorname{Sym}^m A$. The inverse takes the union of the
corresponding horizontal sheets with multiplicity. In flat charts
convergence of such multisets gives Hausdorff convergence of supports
and weak convergence of the integration currents. Over a compact
parameter interval the invariant multisets form a closed subset of
the product of that interval with $\operatorname{Sym}^m A$ and are
therefore compact.

The integral class is locally constant on this family, including its
torsion component. To see this also in the ambient relative cycle
space, take disjoint tubular neighborhoods of the distinct leaf
tori supporting a given cycle. All sufficiently nearby supports lie
in these tubes. In each tube a closed cycle has integral class a
multiple of the core since the tube retracts onto the torus. Pair
with the pullback of an area form of integral one on $B$. Its value
on the core is the positive degree of that leaf. The integral in
each tube is continuous at the given cycle, by a cutoff equal to
one near the part of its support in that tube. Being a multiple of
the core degree, it is locally constant. The argument persists with
the parameter under the smooth identification just constructed.
Thus the locus of each face class is compact and open over the
interval. Its characteristic function yields the endpoint relation
in $S$.

Fix a positive integer $K$. For all sufficiently small $\tau>0$,
every point of $f_\tau$ of period at most $K$ is one of the sixteen
equilibria. Indeed $f_\tau^j x=x$ implies
$\varphi^{2j\tau}x=x$. Write $V$ for the lifted periodic Hamiltonian
vector field on $\RR^4$. Bounded speed ensures that any closed
trajectory on $A$ of sufficiently short period lifts to a closed
trajectory in $\RR^4$. If $V$ has Lipschitz constant $L$ and such a
lift has period $T$, its velocity has zero average. At a point of
maximum speed $M_0$, subtracting this average and bounding the
difference of the two velocities gives $M_0\le LT M_0$. Taking
$2K\tau$ smaller than both short-time bounds excludes every
nonconstant trajectory in question.

At an equilibrium, the return derivative on either plane is minus
the exponential of the linear Hamiltonian flow. A definite Hessian
on that plane gives an elliptic block, and an indefinite Hessian
gives a negative hyperbolic block. We may choose $\tau$ in the
preceding interval to give irrational rotation angles divided by
$\pi$ on all the elliptic blocks, by omitting countably many values.
For the full return derivative $D_v$ at $v/2$ we have
\[
 \operatorname{sign}\det(1-D_v)=+1,\qquad
 \operatorname{sign}\det(1-D_v^2)=(-1)^{\operatorname{ind}_{v/2}(G)} .
\]
Indeed each negative hyperbolic plane reverses precisely the second
sign and has Morse index one. Each elliptic plane has even Morse
index. There are eight odd-index equilibria since each cosine
coordinate independently contributes zero or one negative Hessian
direction. The local return formula of Proposition~\ref{prop:return}
therefore gives $E(Q^{\beta_v})$ at the even-index points and
$E(Q^{\beta_v})/E(Q^{2\beta_v})$ at the odd-index points,
where $E=\Phi(F_1)$. All cycles through degree $K$ at this endpoint
are supported on the sixteen disjoint tori. Excision into disjoint
tubes makes their local series multiply, retaining the integral
label of each summand. The endpoint relation and $2\beta_v=2\beta_0$
now give \eqref{eq:face} after $\Phi$, through degree $K$. Each
cutoff compares with the same complex endpoint, so equality in $S$
holds in every degree.

\smallskip
\noindent\emph{Lifting the identity.}
Let $V_\gamma$ denote the coefficient of either side of
\eqref{eq:face} before applying $\Phi$. The coproduct $\Delta$
induced by addition and pullback $\rho_d$ under multiplication of
cycles give
\[
 \Delta V_\gamma=
    \sum_{\eta+\xi=\gamma}V_\eta\otimes V_\xi,\qquad
 \rho_d V_\gamma=\sum_{d\eta=\gamma}V_\eta .
\]
For the class loci these identities follow from their inverse images.
Nonnegative vertical area puts every effective summand and effective
root of $\gamma$ in $\Gamma^0$. In particular every section label
is a root of the common double. For the product the identities follow
directly on the nontorsion inverse pair defining $F_1$. Its entire cycle
locus of degree $m$ is the point $m[C]$. Under addition its preimage
consists of the ordered pairs of multiples whose degrees sum to $m$.
Under multiplication by $d$ the preimage is empty unless $d\mid m$,
when it is the point $(m/d)[C]$. Substitution $Q^\eta\mapsto Q^{d\eta}$ is a
ring homomorphism even when several monomials have the same image, and
preserves inverses of series with constant term one.

If there were a first base degree of disagreement, the difference
of the coefficients at any such label would be primitive for
$\Delta$, since the summands of smaller positive degree agree.
Every proper division would kill this difference, since all roots
in its division formula have smaller base degree. By the
primitive-division lemma of \cite[Lemma~7.1]{WueProduct} the difference
is a finite linear combination of the multiplicity-one scalar
classes $x_{h,1}$. Their images under $\Phi$ are linearly
independent by \cite[Theorem~5.1]{WueProduct}. The identity in $S$ just
proved makes this difference zero, a contradiction.
\end{proof}

After forgetting the torsion labels, the exponents at degrees one
and two are $16$ and $-8$. Their divisor sums give fixed point
indices $16$ on the odd iterates and $0$ on the even iterates, in
agreement with the Lefschetz numbers of negation on $A$ and the
iteration formula of Dold \cite{Dold}. For comparison, put
$\mathcal E(q)=\prod_{j\ge1}(1-q^j)^{-1}$, the numerical elliptic
partition. The Lefschetz zeta function here is
$\zeta(q)=(1+q)^8/(1-q)^8$, whence
\[
 \mathcal E(q)^{16}\mathcal E(q^2)^{-8}
          =\prod_{j\ge1}\zeta(q^j).
\]
Ionel and Parker use a M\"obius-normalized Gromov series in their
mapping-torus formula. In real dimension six it concerns the
genus-one numerical series. Their distinguished-section hypothesis
requires $\det(1-f_*|_{H_1(A)})=\pm1$, whereas negation on $A$ has
determinant $\det(1-f_*|_{H_1(A)})=16$
\cite[Theorems~0.1 and~4.1]{IPmaps}. Equation~\eqref{eq:face}
retains all sixteen integral section labels before numerical
specialization.
